\section{Modellazione CAD, Fabbricazione e Manutenzione}
%
Il percorso di tirocinio è stato caratterizzato da una forte componente pratica e operativa, che mi ha permesso di consolidare e ampliare significativamente le mie competenze in molteplici ambiti tecnologici. Dal punto di vista tecnico, l'attività si è ramificata in tre aree principali: la progettazione software, la fabbricazione fisica e la manutenzione dell'infrastruttura di laboratorio.

%
\subsection{Modellazione 3D e preparazione file}
L'ambiente dinamico del Fablab e l'eterogeneità dei progetti affrontati hanno richiesto un approccio flessibile alla modellazione CAD, portandomi a utilizzare diversi software a seconda delle specifiche necessità e dell'utenza coinvolta:

\begin{itemize}
    \item \emph{Solidworks}: L'ho utilizzato principalmente per la modellazione e la progettazione meccanica di precisione. Poter applicare in un contesto lavorativo e pratico le nozioni teoriche apprese durante le lezioni è stato fondamentale per consolidare il mio metodo progettuale.
    \item \emph{Fusion360}: È stato lo strumento di elezione per i progetti caratterizzati da forme più organiche o con una spiccata componente artistica, sfruttando la sua flessibilità nella modellazione superficiale.
    \item \emph{OnShape}: Si è rivelato uno strumento strategico e indispensabile durante le attività di tutoring. Essendo basato su cloud, ha permesso agli studenti sprovvisti di licenze o hardware performante di approcciarsi al CAD professionale in modo gratuito e collaborativo.
    \item \emph{Sketchup e AutoCAD}: Sono stati impiegati specificamente per affiancare il gruppo di studenti di architettura. Con questi software abbiamo curato lo studio del terreno, l'analisi delle curve di livello collinari e la modellazione 3D di piante e viste topografiche, successivamente riprodotte in scala.
    \item \emph{Inkscape}: Parallelamente alla progettazione 3D, ho approfondito la grafica vettoriale 2D per la gestione dei tracciati destinati al taglio laser. Questo software è stato essenziale per preparare file puliti, impostare correttamente i livelli di taglio e incisione (ad esempio per la creazione di portachiavi in compensato o pannelli per cabinati).
\end{itemize}

%
\subsection{Reverse engineering}
Non mi sono limitato al solo utilizzo passivo dei macchinari, ma ho curato l'intero ciclo di vita del prodotto. Un aspetto ricorrente è stato il reverse engineering: partendo dall'analisi di componenti fisici rotti o non più in commercio, ho utilizzato strumenti di misurazione per ricavarne le quote, procedendo poi alla modellazione digitale e alla stampa 3D di prototipi iterativi fino al raggiungimento del pezzo funzionale definitivo. Questo approccio è stato spiegato e trasmesso anche agli studenti durante i workshop, illustrando loro il passaggio critico da un'esigenza reale a una produzione di stampo industriale.

%
\subsection{Manutenzione macchinari}
Una parte considerevole e altamente formativa del tirocinio ha riguardato la cura del parco macchine del laboratorio. Ho trovato alcune stampanti 3D non operative e mi sono occupato in prima persona della loro rimessa in funzione. Questo processo ha richiesto un approccio analitico alla risoluzione dei problemi.
Ho analizzato i malfunzionamenti, identificando problemi legati all'estrusione, al blocco meccanico dei motorini delle stampanti e al software.
Ho eseguito poi degli interventi di manutenzione, come lo smontaggio del blocco ``Hot-End'' e dei motorini stepper dedicati al raccoglimento del filamento e ripristinando i relativi cablaggi.